% TikZ 图：伽利略理想斜面实验
% Day 06：牛顿运动定律
\begin{tikzpicture}[scale=0.82,>=Latex]
  % 左图：右侧仍为斜面，小球能上升到与初始位置相同的高度。
  \begin{scope}[shift={(-3.75,0)}]
    \node[font=\footnotesize] at (0,1.95) {(1) 右侧仍为斜面};
    \draw[thick] (-2.8,1.35) -- (0,0) -- (2.8,1.35);
    \draw[dashed,gray] (-2.95,1.35) -- (2.95,1.35);
    \draw[<->,gray] (-3.15,0) -- (-3.15,1.35) node[midway,left] {$h$};

    \filldraw[fill=white,draw=black,thick] (-2.25,1.08) circle (0.13);
    \filldraw[fill=white,draw=black,thick] (2.25,1.08) circle (0.13);
    \draw[->,thick,red] (-1.95,0.93) -- (-1.18,0.56) node[midway,below] {$v$};
    \draw[->,thick,red] (1.15,0.56) -- (1.92,0.93);

    \node[font=\footnotesize] at (0,-0.42) {无摩擦：上升到同一高度};
  \end{scope}

  % 右图：右侧斜面逐渐放平，极限情形为水平面上的匀速直线运动。
  \begin{scope}[shift={(3.75,0)}]
    \node[font=\footnotesize] at (0,1.95) {(2) 右侧放平成水平面};
    \draw[thick] (-2.8,1.35) -- (0,0) -- (3.25,0);
    \draw[dashed,gray] (-2.95,1.35) -- (-1.35,1.35);
    \draw[<->,gray] (-3.15,0) -- (-3.15,1.35) node[midway,left] {$h$};

    \filldraw[fill=white,draw=black,thick] (-2.25,1.08) circle (0.13);
    \filldraw[fill=white,draw=black,thick] (0.75,0.13) circle (0.13);
    \draw[->,thick,red] (-1.95,0.93) -- (-1.18,0.56) node[midway,below] {$v$};
    \draw[->,thick,red] (0.98,0.13) -- (2.28,0.13) node[midway,above] {$v$};

    \node[font=\footnotesize] at (0,-0.42) {无摩擦：速度大小保持};
  \end{scope}
\end{tikzpicture}
